The contemporary understanding of sensory processing in the visual system is founded upon the seminal findings of Hubel and Wiesel (1962, 1968). Under this classical paradigm, the visual cortex is organized as a hierarchy of filters where neurons extract progressively complex features of the visual scene through stage-by-stage processing. Central to this conception is the classical receptive field (CRF), operationally defined as the specific region of visual space where the presence of a light stimulus can evoke a direct neuronal response, significantly influencing its firing rate. In this traditional model, processing has historically been interpreted as a predominantly local phenomenon, where the response of each functional unit is determined almost exclusively by the physical properties of the stimulus located within its defined spatial boundaries.

However, the strictly local view of neuronal processing has been challenged by the discovery of extra-classical receptive field (eCRF) effects. Research has shown that sensory responses are not only determined by stimuli within the CRF but are also significantly modulated by contextual information from the surrounding visual field (Allman et al., 1985). The most prominent of these effects is surround suppression, a non-linear phenomenon where the expansion of a stimulus beyond the CRF limits leads to a substantial decrease in firing rates. This contextual modulation implies a high level of global integration, indicating that the primary visual cortex operates as more than a set of isolated local filters. Such findings suggest that the functional architecture of the visual system must incorporate mechanisms capable of integrating information across disparate spatial locations, a requirement that traditional, purely local hierarchical models fail to satisfy (Angelucci \& Bullier, 2003).

The mechanistic origin of surround suppression remains a central controversy in visual neuroscience. Traditionally, long-range horizontal connections within V1 were proposed as the primary substrate for eCRF effects. However, accumulating evidence suggests that these lateral pathways are physiologically limited by their slow axonal conduction velocities and restricted spatial reach, which fail to account for the rapid and long-range nature of contextual modulation (Bair, 2005; Schwabe et al., 2006). This has led to the prevailing hypothesis that feedback projections from higher-order areas, such as V2, play a decisive role in integrating global information. Top-down signals provide the necessary spatial extent and speed to modulate V1 activity effectively (Angelucci et al., 2017). Reconciling how the laminar microcircuitry reconciles these local lateral and global feedback inputs is essential for a complete model of visual integration.

Despite significant advances in developing biologically plausible cortical microcircuits, such as the canonical Potjans and Diesmann (2014) model, a substantial gap remains between theoretical predictions and experimental observations of eCRF effects. Most existing computational models fail to quantitatively replicate the magnitude of surround suppression observed physiologically, which can reach a reduction of approximately 75\% in neuronal firing rates. This limitation often stems from feedback implementations that are either linear or excessively simplified, thereby failing to capture the intrinsic competitive dynamics of the visual hierarchy. Consequently, there is a critical need for large-scale spiking neural network models that respect laminar architecture to elucidate how inter-areal interactions give rise to robust contextual modulations.

In this study, we propose that robust surround suppression is not merely a byproduct of hierarchical feedback but an emergent property of competitive dynamics within the visual system. Utilizing the NEST simulator, we implemented a large-scale spiking neural network model based on the canonical cortical microcircuit to examine how inter-areal interactions configure eCRF effects. Our central hypothesis holds that lateral competition within higher-order areas, specifically V2, is the essential mechanism for regulating top-down modulation and reaching the suppression levels observed experimentally. By comparing various architectural configurations, this work demonstrates that the integration of inter-areal competition is the critical missing piece for replicating the complex non-linear dynamics of the primary visual cortex.

Beyond its fundamental importance for sensory integration, the characterization of these cortical mechanisms holds significant implications for the emerging field of computational psychiatry. The delicate balance between excitation and inhibition (E/I balance) that governs eCRF effects is also a critical determinant of healthy cognitive function. Disruptions in this equilibrium are hallmark features of several neuropsychiatric disorders, most notably schizophrenia, where patients often exhibit specific deficits in both surround suppression and gamma-band oscillations (Uhlhaas \& Singer, 2015). By providing a platform to systematically manipulate circuit-level parameters, our model allows for the simulation of these pathological endophenotypes. This approach not only validates the biological relevance of our hierarchical competitive framework but also offers a bridge between large-scale circuit dynamics and the search for objective biomarkers in clinical neuroscience.
